% ECONOMICS_PROBLEM: {"id": "D72-389", "title": "Elite capture under inclusive institutions", "jel_code": "D72", "source_id": "OP-0385"}
% !TeX program = xelatex
\documentclass[11pt,a4paper]{article}
\usepackage[margin=23mm]{geometry}
\usepackage{fontspec}
\setmainfont{Latin Modern Roman}
\usepackage{amsmath,amssymb,mathtools}
\usepackage{xcolor,enumitem,booktabs,longtable,array,xurl,hyperref}
\definecolor{muted}{HTML}{53636C}
\hypersetup{colorlinks=true,urlcolor=blue,linkcolor=blue}
\setlength{\parindent}{0pt}
\setlength{\parskip}{5pt}
\newcommand{\R}{\mathbb R}
\newcommand{\E}{\mathbb E}
\newcommand{\N}{\mathbb N}
\newcommand{\ind}{\mathbf 1}
\newcommand{\Var}{\operatorname{Var}}
\newcommand{\Cov}{\operatorname{Cov}}
\newcommand{\supp}{\operatorname{supp}}
\DeclareMathOperator*{\argmin}{arg\,min}
\DeclareMathOperator*{\argmax}{arg\,max}
\newcommand{\entrymeta}[3]{{\sffamily\small\color{muted}\textbf{\texttt{#1}}\quad|\quad #2\par #3\par}}
\newcommand{\Problemref}[1]{\texttt{#1}}
\newcommand{\JELref}[2]{\texttt{#2} (JEL \texttt{#1})}
\begin{document}
\section*{Elite capture under inclusive institutions}
\textbf{Problem ID:} \texttt{D72-389}\quad \textbf{JEL:} \texttt{D72}\par
% BEGIN ECONOMICS STATEMENT
\label{problem:OP-0385}
{\sffamily\small\color{muted}Original subject group: Political economy and institutions\par}
\label{definitions:problem:30007667}
\textbf{Audit classification:} Background; proposed extension. The World Development Report 2017 primary overview and Main Messages are verified background. Their inspected policy framework does not itself state a further-work passage for this exact question; the research target and variants are editorial extensions.
\subsection*{Definitions and precise research target}
\textbf{Complete editorial problem.} Fix a polity undergoing a precisely specified institutional reform \(z\), such as expanded electoral participation or competitive entry into political office. Let \(P_{it}(z)\) be actor \(i\)'s political influence, measured through offices, appointments, policy access, or networks according to a declared rule. Let \(E\) denote actors classified as incumbent elites using prereform information, and define elite influence share \(S_t(z)=\sum_{i\in E}P_{it}(z)/\sum_iP_{it}(z)\), with a positive denominator. Let \(Y_t(z)\) measure broad economic welfare. Estimate changes in \(S\), policy responsiveness, and \(Y\) over a horizon allowing elite adaptation. Decompose descriptive persistence into routes such as party control, financing, social networks, economic ownership, or relatives taking office, while causal channel claims require additional variation. Evaluate complementary reforms \(p\) that enhance contestability and prevent excluded citizens or new elites from recreating capture. Inclusive legal forms do not ensure inclusive policy outcomes. The task adapts the governance report's explicit problem of exclusion and capture; the specific influence metric, dynamic intervention, and replacement-elite test are our specification. Existing country studies offer partial answers, and original priority for this broad formulation is not established.

\textbf{Definitions and admissible scope.} Influence \(P_{it}\ge0\) must use a comparable cardinal measurement, for example normalized voting power or a fixed weighted office-control index; incomparable network and financing scores cannot simply be added. Define prereform elite set \(E\), specify how relatives and controlled organizations are treated, and require \(\sum_iP_{it}>0\). The influence share is descriptive; causal capture means the effect of elite access on decisions relative to a defined access intervention. Define welfare on a fixed citizen population and reform paths that account for elite adaptation, replacement elites, and coalition changes. Legal enfranchisement, effective participation, and policy responsiveness are separate measured objects. Fix a finite actor universe including potential entrants, \(T<\infty\), a prereform elite-classification rule, and nonnegative comparable influence scores. Removed actors receive score zero, and the total score must be positive under every regime. Actor descendants and controlled organizations are assigned by the declared lineage rule without duplicate attribution. Let \(\tau_S(T,z)=E[S_T(z)-S_T(0)]\), with welfare contrast \(\tau_Y\) defined on the fixed citizen population and an explicit welfare scale. These are policy effects, whereas channel contributions require additional interventions. The target includes feasible complement policies and their joint effects under the declared assignment, influence-measurement, and equilibrium-response law.
\subsection*{Variants and assumption changes}
\begin{enumerate}
\item Incumbent lineage persistence (editorial informed by governance agenda): track original elites, their relatives, and controlled parties separately through term limits or franchise expansion; compare office turnover with effective policy-access change.
\item Contestability bundle (source-motivated extension): combine participation reform with financing disclosure or competitive candidate entry, explicitly defining enforcement. Estimate whether policy responsiveness shifts toward previously excluded citizens while monitoring new elite concentration.
\end{enumerate}
\subsection*{References and source audit}
\textbf{Support and access.} Primary report page checked. Exclusion, capture and contestability motivate the agenda; the quantitative elite-share metric and lineage tests are editorial. \textbf{Locator:} WDR 2017 About, paragraphs 1--3 and Main Messages. \emph{Provenance limit.} The inspected overview presents the report's questions, policy framework, and conclusions. It does not establish the precise editorial problem here as an explicitly published research agenda or open conjecture.
\begin{enumerate}\raggedright
\item\hypertarget{definitions:source:30007667:1}{} World Bank. 2017. \emph{World Development Report 2017: Governance and the Law}. About, opening three paragraphs; Main Messages, first three bullets.
\url{https://www.worldbank.org/en/publication/wdr2017}.
\end{enumerate}

\subsection*{Common conventions: Source mathematical conventions}

These conventions apply to each entry unless explicitly replaced.

\textbf{Models and identification.} A primitive model \(m\) specifies all probability laws, feasible choices, information, equilibrium and measurement rules used by its target. The admissible class \(\mathcal M\) is fixed independently of outcomes. If \(P_O(m)\) is its observable probability law and \(\tau(m)\) the target, its identified set at observable law \(P\) is
\[
 \mathcal I_\tau(P)=\{\tau(m):m\in\mathcal M,\ P_O(m)=P\}.
\]
Point identification means that this set is a singleton. An empty set signals incompatibility of data and assumptions. Coordinate bounds need not describe the joint set. Bounds are sharp only if every retained target is attainable or arbitrarily approachable by compatible models. Finite bounds require explicit restrictions on unobserved quantities; unrestricted confounding, hidden wealth, extrapolation or arbitrary equilibrium selection can yield uninformative sets.

\textbf{Causal interventions.} A potential outcome \(Y_i(p)\) is the outcome under a complete policy regime \(p\), including timing, financing, information and market spillovers. The target population and units are fixed before treatment. With interference, use the complete assignment vector or a justified exposure mapping; individual no-interference notation alone cannot identify an economy-wide intervention. A difference of mean potential outcomes does not require the cross-world joint distribution. Conditioning on post-treatment migration, survival, enrollment or employment changes the estimand and needs separate justification.

\textbf{Statistical statements.} An estimator, confidence set, forecast or test specifies a sampling law, model class and loss/coverage criterion. Uniform \((1-\alpha)\)-coverage means \(\inf_{m\in\mathcal M}\Pr_m\{\tau(m)\in C_n\}\geq1-\alpha\), or an explicitly asymptotic version. The whole parameter space trivially has coverage, so informativeness requires a stated width, risk or power objective. A forecasting improvement is relative to a named benchmark, information set and evaluation population.

\textbf{Equilibrium and optimization.} An equilibrium comprises feasible plans, optimality, beliefs where needed, budget constraints and all market/resource clearing. The primitives or a stated benchmark define each correspondence \(\mathcal E(p;m)\). Nonemptiness is assumed or proved, never inferred just from compact policy sets. A selected-equilibrium target uses a declared selection rule; a robust target quantifies over all permitted equilibria. Use suprema or infima unless attainment is established. An \(\varepsilon\)-optimal policy is feasible and within \(\varepsilon\) of the finite optimum value. Report an empty feasible set separately.

\textbf{Welfare and accounting.} Normative utility, interpersonal normalization, social weights, numeraire, discounting and the use of public receipts are primitives. Behavioral utility need not coincide with the welfare criterion. Household welfare includes ownership income and rents once; adding profits again to compensating variation generally double-counts them. A monetary equivalent is defined by an explicit transfer and price convention, with monotonicity and range conditions. Average effects, mechanism contributions, transfers and welfare losses are different targets. Infinite sums require convergence and dynamic feasible plans require terminal or no-Ponzi conditions.

\textbf{Source versions.} A working-paper date, revision date and journal publication year are distinguished. A DOI for a journal version is not automatically the DOI of an earlier manuscript. Locators refer to the stated version and printed pages where specified. A metadata-only check does not verify a claim about a full-text conclusion. Every entry's source subsection narrows the provenance claim accordingly.
% END ECONOMICS STATEMENT
\end{document}
